% Options for packages loaded elsewhere
\PassOptionsToPackage{unicode}{hyperref}
\PassOptionsToPackage{hyphens}{url}
\documentclass[
]{article}
\usepackage{xcolor}
\usepackage[margin=1in]{geometry}
\usepackage{amsmath,amssymb}
\setcounter{secnumdepth}{-\maxdimen} % remove section numbering
\usepackage{iftex}
\ifPDFTeX
  \usepackage[T1]{fontenc}
  \usepackage[utf8]{inputenc}
  \usepackage{textcomp} % provide euro and other symbols
\else % if luatex or xetex
  \usepackage{unicode-math} % this also loads fontspec
  \defaultfontfeatures{Scale=MatchLowercase}
  \defaultfontfeatures[\rmfamily]{Ligatures=TeX,Scale=1}
\fi
\usepackage{lmodern}
\ifPDFTeX\else
  % xetex/luatex font selection
\fi
% Use upquote if available, for straight quotes in verbatim environments
\IfFileExists{upquote.sty}{\usepackage{upquote}}{}
\IfFileExists{microtype.sty}{% use microtype if available
  \usepackage[]{microtype}
  \UseMicrotypeSet[protrusion]{basicmath} % disable protrusion for tt fonts
}{}
\makeatletter
\@ifundefined{KOMAClassName}{% if non-KOMA class
  \IfFileExists{parskip.sty}{%
    \usepackage{parskip}
  }{% else
    \setlength{\parindent}{0pt}
    \setlength{\parskip}{6pt plus 2pt minus 1pt}}
}{% if KOMA class
  \KOMAoptions{parskip=half}}
\makeatother
\usepackage{color}
\usepackage{fancyvrb}
\newcommand{\VerbBar}{|}
\newcommand{\VERB}{\Verb[commandchars=\\\{\}]}
\DefineVerbatimEnvironment{Highlighting}{Verbatim}{commandchars=\\\{\}}
% Add ',fontsize=\small' for more characters per line
\usepackage{framed}
\definecolor{shadecolor}{RGB}{248,248,248}
\newenvironment{Shaded}{\begin{snugshade}}{\end{snugshade}}
\newcommand{\AlertTok}[1]{\textcolor[rgb]{0.94,0.16,0.16}{#1}}
\newcommand{\AnnotationTok}[1]{\textcolor[rgb]{0.56,0.35,0.01}{\textbf{\textit{#1}}}}
\newcommand{\AttributeTok}[1]{\textcolor[rgb]{0.13,0.29,0.53}{#1}}
\newcommand{\BaseNTok}[1]{\textcolor[rgb]{0.00,0.00,0.81}{#1}}
\newcommand{\BuiltInTok}[1]{#1}
\newcommand{\CharTok}[1]{\textcolor[rgb]{0.31,0.60,0.02}{#1}}
\newcommand{\CommentTok}[1]{\textcolor[rgb]{0.56,0.35,0.01}{\textit{#1}}}
\newcommand{\CommentVarTok}[1]{\textcolor[rgb]{0.56,0.35,0.01}{\textbf{\textit{#1}}}}
\newcommand{\ConstantTok}[1]{\textcolor[rgb]{0.56,0.35,0.01}{#1}}
\newcommand{\ControlFlowTok}[1]{\textcolor[rgb]{0.13,0.29,0.53}{\textbf{#1}}}
\newcommand{\DataTypeTok}[1]{\textcolor[rgb]{0.13,0.29,0.53}{#1}}
\newcommand{\DecValTok}[1]{\textcolor[rgb]{0.00,0.00,0.81}{#1}}
\newcommand{\DocumentationTok}[1]{\textcolor[rgb]{0.56,0.35,0.01}{\textbf{\textit{#1}}}}
\newcommand{\ErrorTok}[1]{\textcolor[rgb]{0.64,0.00,0.00}{\textbf{#1}}}
\newcommand{\ExtensionTok}[1]{#1}
\newcommand{\FloatTok}[1]{\textcolor[rgb]{0.00,0.00,0.81}{#1}}
\newcommand{\FunctionTok}[1]{\textcolor[rgb]{0.13,0.29,0.53}{\textbf{#1}}}
\newcommand{\ImportTok}[1]{#1}
\newcommand{\InformationTok}[1]{\textcolor[rgb]{0.56,0.35,0.01}{\textbf{\textit{#1}}}}
\newcommand{\KeywordTok}[1]{\textcolor[rgb]{0.13,0.29,0.53}{\textbf{#1}}}
\newcommand{\NormalTok}[1]{#1}
\newcommand{\OperatorTok}[1]{\textcolor[rgb]{0.81,0.36,0.00}{\textbf{#1}}}
\newcommand{\OtherTok}[1]{\textcolor[rgb]{0.56,0.35,0.01}{#1}}
\newcommand{\PreprocessorTok}[1]{\textcolor[rgb]{0.56,0.35,0.01}{\textit{#1}}}
\newcommand{\RegionMarkerTok}[1]{#1}
\newcommand{\SpecialCharTok}[1]{\textcolor[rgb]{0.81,0.36,0.00}{\textbf{#1}}}
\newcommand{\SpecialStringTok}[1]{\textcolor[rgb]{0.31,0.60,0.02}{#1}}
\newcommand{\StringTok}[1]{\textcolor[rgb]{0.31,0.60,0.02}{#1}}
\newcommand{\VariableTok}[1]{\textcolor[rgb]{0.00,0.00,0.00}{#1}}
\newcommand{\VerbatimStringTok}[1]{\textcolor[rgb]{0.31,0.60,0.02}{#1}}
\newcommand{\WarningTok}[1]{\textcolor[rgb]{0.56,0.35,0.01}{\textbf{\textit{#1}}}}
\usepackage{graphicx}
\makeatletter
\newsavebox\pandoc@box
\newcommand*\pandocbounded[1]{% scales image to fit in text height/width
  \sbox\pandoc@box{#1}%
  \Gscale@div\@tempa{\textheight}{\dimexpr\ht\pandoc@box+\dp\pandoc@box\relax}%
  \Gscale@div\@tempb{\linewidth}{\wd\pandoc@box}%
  \ifdim\@tempb\p@<\@tempa\p@\let\@tempa\@tempb\fi% select the smaller of both
  \ifdim\@tempa\p@<\p@\scalebox{\@tempa}{\usebox\pandoc@box}%
  \else\usebox{\pandoc@box}%
  \fi%
}
% Set default figure placement to htbp
\def\fps@figure{htbp}
\makeatother
\setlength{\emergencystretch}{3em} % prevent overfull lines
\providecommand{\tightlist}{%
  \setlength{\itemsep}{0pt}\setlength{\parskip}{0pt}}
\usepackage{bookmark}
\IfFileExists{xurl.sty}{\usepackage{xurl}}{} % add URL line breaks if available
\urlstyle{same}
\hypersetup{
  pdftitle={ANCOVA in RCTs (baseline adjustment)},
  hidelinks,
  pdfcreator={LaTeX via pandoc}}

\title{ANCOVA in RCTs (baseline adjustment)}
\author{}
\date{\vspace{-2.5em}}

\begin{document}
\maketitle

\textbf{Inference labs} use the five-step template: Hypothesis →
Visualise → Assumptions → Conduct → Conclude.

\subsection{Learning objectives}\label{learning-objectives}

\begin{itemize}
\tightlist
\item
  Choose between change scores, post-only, and ANCOVA in an RCT analysis
  plan.
\item
  Fit an ANCOVA model and read the treatment effect as an adjusted-mean
  difference.
\item
  Show that ANCOVA is more efficient than the change-score approach when
  baseline and follow-up are correlated.
\end{itemize}

\subsection{Prerequisites}\label{prerequisites}

Linear regression; one-way ANOVA.

\subsection{Background}\label{background}

In a randomised trial with a continuous outcome measured at baseline and
follow-up, three analyses are common: compare the post-intervention
means, compare the change scores, or use analysis of covariance (ANCOVA)
with baseline as a covariate. Randomisation ensures all three are
unbiased in expectation, but they are not equally efficient.

ANCOVA weights the baseline by its observed correlation with the
follow-up and is therefore always at least as efficient as the
change-score approach, and much more efficient when that correlation is
moderate. Change scores are equivalent to ANCOVA with the baseline slope
constrained to 1; post-only is ANCOVA with the slope constrained to 0.
When neither constraint matches the data, efficiency is lost.

Regression to the mean is the mechanical reason ANCOVA wins. Patients
with unusually high baselines tend to have lower follow-ups whether they
are treated or not; ANCOVA uses that pattern, change-score analysis
pretends it does not exist.

\subsection{Setup}\label{setup}

\begin{Shaded}
\begin{Highlighting}[]
\FunctionTok{library}\NormalTok{(tidyverse)}
\FunctionTok{library}\NormalTok{(broom)}
\FunctionTok{set.seed}\NormalTok{(}\DecValTok{42}\NormalTok{)}
\FunctionTok{theme\_set}\NormalTok{(}\FunctionTok{theme\_minimal}\NormalTok{(}\AttributeTok{base\_size =} \DecValTok{12}\NormalTok{))}
\end{Highlighting}
\end{Shaded}

\subsection{1. Hypothesis}\label{hypothesis}

Simulate a trial with baseline and follow-up correlated at r = 0.6, then
compare the three analyses.

\subsection{2. Visualise}\label{visualise}

\begin{Shaded}
\begin{Highlighting}[]
\NormalTok{baseline }\OtherTok{\textless{}{-}} \FunctionTok{rnorm}\NormalTok{(n, }\DecValTok{140}\NormalTok{, }\DecValTok{15}\NormalTok{)}
\NormalTok{arm }\OtherTok{\textless{}{-}} \FunctionTok{rep}\NormalTok{(}\FunctionTok{c}\NormalTok{(}\StringTok{"placebo"}\NormalTok{, }\StringTok{"active"}\NormalTok{), }\AttributeTok{each =}\NormalTok{ n }\SpecialCharTok{/} \DecValTok{2}\NormalTok{)}
\CommentTok{\# true treatment effect on follow{-}up: {-}5}
\NormalTok{followup }\OtherTok{\textless{}{-}} \FloatTok{0.6} \SpecialCharTok{*}\NormalTok{ (baseline }\SpecialCharTok{{-}} \DecValTok{140}\NormalTok{) }\SpecialCharTok{+} \DecValTok{140} \SpecialCharTok{+}
  \FunctionTok{ifelse}\NormalTok{(arm }\SpecialCharTok{==} \StringTok{"active"}\NormalTok{, }\SpecialCharTok{{-}}\DecValTok{5}\NormalTok{, }\DecValTok{0}\NormalTok{) }\SpecialCharTok{+} \FunctionTok{rnorm}\NormalTok{(n, }\DecValTok{0}\NormalTok{, }\DecValTok{12}\NormalTok{)}
\NormalTok{dat }\OtherTok{\textless{}{-}} \FunctionTok{tibble}\NormalTok{(baseline, followup, arm,}
              \AttributeTok{change =}\NormalTok{ followup }\SpecialCharTok{{-}}\NormalTok{ baseline) }\SpecialCharTok{|\textgreater{}}
  \FunctionTok{mutate}\NormalTok{(}\AttributeTok{arm =} \FunctionTok{factor}\NormalTok{(arm, }\AttributeTok{levels =} \FunctionTok{c}\NormalTok{(}\StringTok{"placebo"}\NormalTok{, }\StringTok{"active"}\NormalTok{)))}

\FunctionTok{ggplot}\NormalTok{(dat, }\FunctionTok{aes}\NormalTok{(baseline, followup, }\AttributeTok{colour =}\NormalTok{ arm)) }\SpecialCharTok{+}
  \FunctionTok{geom\_point}\NormalTok{(}\AttributeTok{alpha =} \FloatTok{0.6}\NormalTok{) }\SpecialCharTok{+}
  \FunctionTok{geom\_smooth}\NormalTok{(}\AttributeTok{method =} \StringTok{"lm"}\NormalTok{, }\AttributeTok{se =} \ConstantTok{FALSE}\NormalTok{) }\SpecialCharTok{+}
  \FunctionTok{labs}\NormalTok{(}\AttributeTok{x =} \StringTok{"Baseline"}\NormalTok{, }\AttributeTok{y =} \StringTok{"Follow{-}up"}\NormalTok{)}
\end{Highlighting}
\end{Shaded}

\subsection{3. Assumptions}\label{assumptions}

Randomisation; linear relationship between baseline and follow-up; equal
slopes across arms (tested with an interaction term).

\begin{Shaded}
\begin{Highlighting}[]
\FunctionTok{tidy}\NormalTok{(fit\_int) }\SpecialCharTok{|\textgreater{}} \FunctionTok{filter}\NormalTok{(term }\SpecialCharTok{==} \StringTok{"armactive:baseline"}\NormalTok{)}
\end{Highlighting}
\end{Shaded}

If the interaction is small, pool the slopes in the ANCOVA.

\subsection{4. Conduct}\label{conduct}

\begin{Shaded}
\begin{Highlighting}[]
\NormalTok{fit\_change }\OtherTok{\textless{}{-}} \FunctionTok{lm}\NormalTok{(change }\SpecialCharTok{\textasciitilde{}}\NormalTok{ arm, }\AttributeTok{data =}\NormalTok{ dat)}
\NormalTok{fit\_ancova }\OtherTok{\textless{}{-}} \FunctionTok{lm}\NormalTok{(followup }\SpecialCharTok{\textasciitilde{}}\NormalTok{ arm }\SpecialCharTok{+}\NormalTok{ baseline, }\AttributeTok{data =}\NormalTok{ dat)}

\FunctionTok{bind\_rows}\NormalTok{(}
  \FunctionTok{tidy}\NormalTok{(fit\_post, }\AttributeTok{conf.int =} \ConstantTok{TRUE}\NormalTok{) }\SpecialCharTok{|\textgreater{}} \FunctionTok{mutate}\NormalTok{(}\AttributeTok{model =} \StringTok{"post only"}\NormalTok{),}
  \FunctionTok{tidy}\NormalTok{(fit\_change, }\AttributeTok{conf.int =} \ConstantTok{TRUE}\NormalTok{) }\SpecialCharTok{|\textgreater{}} \FunctionTok{mutate}\NormalTok{(}\AttributeTok{model =} \StringTok{"change"}\NormalTok{),}
  \FunctionTok{tidy}\NormalTok{(fit\_ancova, }\AttributeTok{conf.int =} \ConstantTok{TRUE}\NormalTok{) }\SpecialCharTok{|\textgreater{}} \FunctionTok{mutate}\NormalTok{(}\AttributeTok{model =} \StringTok{"ANCOVA"}\NormalTok{)}
\NormalTok{) }\SpecialCharTok{|\textgreater{}}
  \FunctionTok{filter}\NormalTok{(term }\SpecialCharTok{==} \StringTok{"armactive"}\NormalTok{) }\SpecialCharTok{|\textgreater{}}
\NormalTok{  dplyr}\SpecialCharTok{::}\FunctionTok{select}\NormalTok{(model, estimate, std.error, conf.low, conf.high)}
\end{Highlighting}
\end{Shaded}

The three estimates should all be around −5; the ANCOVA standard error
should be the smallest.

\subsection{5. Concluding statement}\label{concluding-statement}

\begin{quote}
In a simulated trial (\emph{n} = \texttt{nrow(dat)}), ANCOVA estimated
the treatment effect at
\texttt{round(coef(fit\_ancova){[}"armactive"{]},\ 2)} units (95\% CI
\texttt{round(confint(fit\_ancova){[}"armactive",\ 1{]},\ 2)} to
\texttt{round(confint(fit\_ancova){[}"armactive",\ 2{]},\ 2)}) with a
smaller standard error than either the post-only or change-score
analyses. ANCOVA is the prespecified primary analysis in most
continuous-outcome RCTs.
\end{quote}

\subsection{Common pitfalls}\label{common-pitfalls}

\begin{itemize}
\tightlist
\item
  Adjusting for post-baseline covariates (not covered here, but a
  classic trap).
\item
  Using a change score when baseline and follow-up are strongly
  correlated.
\item
  Dropping observations with missing baselines instead of imputing or
  prespecifying a policy.
\end{itemize}

\subsection{Further reading}\label{further-reading}

\begin{itemize}
\tightlist
\item
  Senn S (2006), \emph{Change from baseline and analysis of
  covariance\ldots{}}
\item
  Vickers AJ, Altman DG (2001), \emph{Analysing controlled trials
  with\ldots{}}
\item
  ICH E9 (R1) Addendum on Estimands.
\end{itemize}

\subsection{Session info}\label{session-info}

\begin{Shaded}
\begin{Highlighting}[]

\end{Highlighting}
\end{Shaded}

\subsection{See also --- chapter index}\label{see-also-chapter-index}

\begin{itemize}
\tightlist
\item
  \href{../index.Rmd}{Methods in this chapter}
\item
  \href{../../../ch00-find-your-method.Rmd}{Find your method (Chapter
  0)}
\end{itemize}

\end{document}
